\addtocontents{toc}{\protect\setcounter{tocdepth}{0}}
\chapter*{LỜI MỞ ĐẦU}\label{loimodau_intro}
\addcontentsline{toc}{chapter}{LỜI MỞ ĐẦU}
\vspace{-40pt}

\renewcommand{\thesection}{\arabic{section}} % Loại bỏ số chương (0) phía trước, chỉ hiện số thứ tự section
% ------------------------------------------
\section{Lý do chọn đề tài}
Máy tính và linh kiện máy tính là những sản phẩm có nhiều chủng loại, thông số kỹ thuật và mức giá khác nhau. Khi hoạt động kinh doanh chủ yếu dựa trên hình thức trực tiếp, khách hàng có thể mất thời gian tìm hiểu sản phẩm, còn cửa hàng gặp khó khăn trong việc giới thiệu thông tin và hỗ trợ quá trình đặt mua. Vì vậy, việc ứng dụng công nghệ thông tin để xây dựng một kênh bán hàng trực tuyến là cần thiết.

Xuất phát từ nhu cầu đó, nhóm lựa chọn đề tài ``Xây dựng website thương mại điện tử bán máy tính và linh kiện máy tính'' cho T\&T Computer. Đề tài hướng đến việc xây dựng một website giúp khách hàng tra cứu sản phẩm, quản lý giỏ hàng và đặt hàng thuận tiện; đồng thời tạo nền tảng để cửa hàng quản lý một số hoạt động bán hàng trên hệ thống. Đây cũng là cơ hội để nhóm vận dụng kiến thức về phân tích, thiết kế cơ sở dữ liệu và phát triển ứng dụng web vào một bài toán cụ thể.
\section{Mục tiêu}
Mục tiêu của đồ án là phân tích, thiết kế và xây dựng website thương mại điện tử T\&T Computer phục vụ việc giới thiệu và bán máy tính, linh kiện máy tính. Website cần cung cấp thông tin sản phẩm và hỗ trợ khách hàng thực hiện các bước mua hàng trực tuyến cơ bản.

Cụ thể, đồ án hướng đến các mục tiêu sau:
\begin{itemize}
    \item Vận dụng kiến thức các môn học liên quan để phân tích yêu cầu, thiết kế cơ sở dữ liệu và xây dựng ứng dụng web.
    \item Xây dựng các chức năng xem danh mục, tìm kiếm và tra cứu thông tin sản phẩm.
    \item Hỗ trợ khách hàng đăng ký, đăng nhập, quản lý hồ sơ, thêm sản phẩm vào giỏ hàng và đặt hàng theo hình thức thanh toán khi nhận hàng.
    \item Hỗ trợ khách hàng theo dõi đơn hàng và gửi yêu cầu hủy đơn theo trạng thái phù hợp; xây dựng chức năng quản trị tài khoản nhân viên.
    \item Tổ chức mã nguồn và cơ sở dữ liệu để các chức năng của website có thể phối hợp trong quy trình bán hàng.
\end{itemize}
\section{Đối tượng và phạm vi}
\begin{itemize}
    \item \textbf{Đối tượng}
    
    Quy trình bán máy tính và linh kiện máy tính trực tuyến; các chức năng của website thương mại điện tử dành cho khách hàng và hoạt động quản trị tài khoản nhân viên; cùng các kỹ thuật xây dựng ứng dụng web và tổ chức cơ sở dữ liệu liên quan.
    \item \textbf{Phạm vi}
    
    Đồ án tập trung xây dựng website T\&T Computer với các chức năng tra cứu sản phẩm, quản lý tài khoản khách hàng, giỏ hàng, đặt hàng COD hoặc thanh toán trực tuyến qua Cổng thanh toán, theo dõi và hủy đơn hàng, cùng quản trị tài khoản nhân viên. Hệ thống được phát triển bằng PHP, sử dụng MySQL để lưu trữ dữ liệu và PDO để kết nối cơ sở dữ liệu; giao diện sử dụng HTML, CSS và JavaScript.
\end{itemize}


\section{Phương pháp}
\begin{itemize}
    \item \textbf{Phương pháp nghiên cứu tài liệu:} Tìm hiểu kiến thức về website thương mại điện tử, quy trình bán hàng trực tuyến, PHP, MySQL và các công nghệ liên quan để làm cơ sở xây dựng hệ thống.
    \item \textbf{Phương pháp phân tích và thiết kế:} Khảo sát yêu cầu của bài toán, xác định các chức năng cần có, phân tích quy trình nghiệp vụ và thiết kế cơ sở dữ liệu phù hợp.
    \item \textbf{Phương pháp xây dựng và kiểm thử:} Thực hiện website bằng PHP, HTML, CSS và JavaScript; kết nối cơ sở dữ liệu MySQL thông qua PDO; kiểm tra các chức năng theo những luồng sử dụng chính như đăng nhập, quản lý giỏ hàng, đặt hàng và xử lý đơn hàng.
\end{itemize}


\section{Cấu trúc tài liệu đồ án}
\begin{itemize}
    \item \textbf{Chương 1: Cơ sở lý thuyết và công nghệ} \\
    Trình bày tổng quan về đề tài, các kiến thức và công nghệ liên quan đến website thương mại điện tử làm cơ sở cho quá trình xây dựng hệ thống.
    
    \item \textbf{Chương 2: Phân tích và thiết kế hệ thống} \\
    Trình bày phương pháp thực hiện, kết quả phân tích yêu cầu, thiết kế hệ thống và cơ sở dữ liệu, cùng các công cụ được sử dụng.
    
    \item \textbf{Chương 3: Phân tích chức năng} \\
    Trình bày các chức năng chính đã xây dựng và cách thức triển khai, cùng các luồng xử lý nghiệp vụ.
    
    \item \textbf{Chương 4: Tổng kết} \\
    Tổng kết kết quả đạt được, nêu những hạn chế và đề xuất hướng phát triển tiếp theo.
\end{itemize}